\hypertarget{openbmc-host-console-support}{%
\section{OpenBMC host console
support}\label{openbmc-host-console-support}}

This document describes how to connect to the host UART console from an
OpenBMC management system.

The console infrastructure allows multiple shared connections to a
single host UART. UART data from the host is output to all connections,
and input from any connection is sent to the host.

\hypertarget{remote-console-connections}{%
\subsection{Remote console
connections}\label{remote-console-connections}}

To connect to an OpenBMC console session remotely, just ssh to your BMC
on port 2200. Use the same login credentials you would for a normal ssh
session:

\begin{verbatim}
ssh -p 2200 [user]@[bmc-hostname]
\end{verbatim}

\hypertarget{local-console-connections}{%
\subsection{Local console connections}\label{local-console-connections}}

If you're already logged into an OpenBMC machine, you can start a
console session directly, using:

\begin{verbatim}
obmc-console-client
\end{verbatim}

To exit from a console, type:

\begin{verbatim}
    return ~ .
\end{verbatim}

Note that if you're on an ssh connection, you'll need to `escape' the
\textasciitilde{} character, by entering it twice:

\begin{verbatim}
    return ~ ~ .
\end{verbatim}

This is because obmc-console-client is an ssh session, and a double
\texttt{\textasciitilde{}} is required to escape the ``inner''
(obmc-console-client) ssh session.

\hypertarget{logging}{%
\subsection{Logging}\label{logging}}

Console logs are kept in: \texttt{/var/log/obmc-console.log}

This log is limited in size, and will wrap after hitting that limit
(currently set at 16kB).
